\documentclass{report}
\usepackage{amsmath,amssymb}
\setlength{\parindent}{0mm}
\setlength{\parskip}{1em}
\begin{document}
\begin{center}
\rule{6in}{1pt} \
{\large Jonathan Hu \\
{\bf MueLu: A Flexible, Parallel Multigrid Framework}}

PO Box 969 \\ MS 9159 \\ Livermore \\ CA 94551-0969
\\
{\tt jhu@sandia.gov}\\
Jeremie Gaidamour\\
Ray Tuminaro\\
Tobias Wiesner\\
Axel Gerstenberger\\
Chris Siefert\end{center}

We have developed a flexible parallel multigrid framework called MueLu
for developing and deploying multigrid algorithms.
The framework is designed with a user-friendly front end, and so provides
the standard complement of smoothers and direct solvers.
Prolongator algorithms include smoothed aggregation and Petrov Galerkin
methods, with plans to include others such as energy
minimization.
The design allows advanced users to customize almost any aspect of the
multigrid solver, e.g., introducing new smoothers,
alternative coarsening algorithms, specialized strength-of-connection
measures, and even hybrids of different multigrid
algorithms. MueLu leverages existing capabilities from Sandia's Trilinos
project in order to manage issues such as parallelism and data types. In
this talk we'll discuss the main design features, illustrate library
usage with examples, and provide parallel examples.


\end{document}
